%%%%%%%%%%%%%%%%%%%%%%%%%%%%%%%%%%%%%%%%%%%%%%%%%%%%%%%%%%%%%%%%%%%%%%%%%%
%%%%%%%%%%%%%%%%%%%%%%%%%%   Paraboloide N Cone    %%%%%%%%%%%%%%%%%%%%%%%%
%%%%%%%%%%%%%%%%%%%%%%%%%%%%%%%%%%%%%%%%%%%%%%%%%%%%%%%%%%%%%%%%%%%%%%%%%%

% Corte, num plano coordenado, do paraboloide v = h^2 e do cone v = |h|:
% a parábola, as duas semirretas, os pontos de interseção (0,0) e (±1,1)
% e parte da região entre as curvas (h^2 <= v <= |h|) sombreada.
% Uso: \ParaboloideNcone{eixo horizontal}{eixo vertical}{q}
% Ex.: \ParaboloideNcone{y}{z}{o}  % a região toda, no plano yz
% Como a região fica toda acima do eixo horizontal, q tem só estas opções:
% q = o, t       : a região toda
% q = 1, d, r    : a parte da direita (1º quadrante)
% q = 2, e, l    : a parte da esquerda (2º quadrante)

\makeatletter

% Comando auxiliar, com a região definida explicitamente
% #1, #2: eixos (horizontal e vertical)   #3, #4: domínio do sombreado, de #3 a #4
\newcommand{\ParaboloideNconeAux}[4]{%
\begin{tikzpicture}[scale=1]
  \begin{axis}[
      xlabel={$#1$}, ylabel={$#2$},
      xtick={-1,0,1}, ytick={0,1},
      x tick label style={anchor=north east}, % rótulos abaixo e à esquerda do eixo horizontal
      y tick label style={anchor=south east}, % rótulos acima e à esquerda do eixo vertical
      width=8cm, height=6cm,
      every axis x label/.style={at={(ticklabel* cs:0.97)}, anchor=north},
      every axis y label/.style={at={(ticklabel* cs:0.97)}, anchor=east},
    ]

    % --- Linhas tracejadas ---
    \addplot[Tracejado] coordinates {(-1,0) (-1,1)};
    \addplot[Tracejado] coordinates {(-1,1) (1,1)};
    \addplot[Tracejado] coordinates {(1,0) (1,1)};

    % --- Curvas ---
    \addplot[name path=parabola, domain=-1.2:1.2, samples=100, thick, blue] {x^2}
      node[pos=0.7, below right, blue] {$#2 = #1^2$};
    \addplot[name path=modulo, domain=-1.2:1.2, samples=100, thick, red] {abs(x)}
      node[pos=0.8, above left, red] {$#2 = |#1|$};

    % --- Região sombreada entre as curvas, no domínio escolhido (q) ---
    \addplot[gray!30] fill between[of=parabola and modulo, soft clip={domain=#3:#4}];

    % --- Pontos de interseção ---
    \addplot[only marks, black!60, mark=*] coordinates {(-1,1) (1,1) (0,0)};

  \end{axis}
\end{tikzpicture}%
}

% Tabela de q -> domínio. Cada entrada recebe #1, #2 (eixos).
\@namedef{ParaboloideNcone@o}#1#2{\ParaboloideNconeAux{#1}{#2}{-1}{1}}
\@namedef{ParaboloideNcone@t}#1#2{\ParaboloideNconeAux{#1}{#2}{-1}{1}}
\@namedef{ParaboloideNcone@1}#1#2{\ParaboloideNconeAux{#1}{#2}{0}{1}}
\@namedef{ParaboloideNcone@d}#1#2{\ParaboloideNconeAux{#1}{#2}{0}{1}}
\@namedef{ParaboloideNcone@r}#1#2{\ParaboloideNconeAux{#1}{#2}{0}{1}}
\@namedef{ParaboloideNcone@2}#1#2{\ParaboloideNconeAux{#1}{#2}{-1}{0}}
\@namedef{ParaboloideNcone@e}#1#2{\ParaboloideNconeAux{#1}{#2}{-1}{0}}
\@namedef{ParaboloideNcone@l}#1#2{\ParaboloideNconeAux{#1}{#2}{-1}{0}}

\newcommand{\ParaboloideNcone}[3]{%
  \@ifundefined{ParaboloideNcone@#3}%
    {\PackageError{ParaboloideNcone}{q = '#3' invalido}{Use o, t, 1, d, r, 2, e ou l (3, 4 e b ficam abaixo do eixo, onde nao ha regiao).}}%
    {\@nameuse{ParaboloideNcone@#3}{#1}{#2}}%
}
\makeatother
